\chapter{2010年北京卷（理）}

\section{单选题}

\begin{problem}
集合 \(P = \{x \in \Z \mid 0 \leqslant x < 3\}\)，\(M = \{x \in \Z \mid x^2 \leqslant 9\}\)，则 \(P \cap M =\)（\quad）

\choices
    {\( \{1, 2\} \)}
    {\(\{0, 1, 2\}\)}
    {\(\{x\mid 0\leqslant x < 3\}\)}
    {\(\{x\mid 0\leqslant x\leqslant 3\}\)}
\end{problem}

\begin{answer}
B
\end{answer}

\begin{solution}
由 \(0\leqslant x<3\) 且 \(x\in\Z\)，得 \(P=\{0,1,2\}\). 又由 \(x^{2}\leqslant9\) 且 \(x\in\Z\)，得 \(M=\{-3,-2,-1,0,1,2,3\}\). 因此 \(P\cap M=\{0,1,2\}\)，故选 B.
\end{solution}

\begin{problem}
在等比数列 \(\{a_{n}\}\) 中，\(a_{1} = 1\)，公比 \(|q| \neq 1\). 若 \(a_{m} = a_{1}a_{2}a_{3}a_{4}a_{5}\)，则 \(m =\)（\quad）

\choices
    {\(9\)}
    {\(10\)}
    {\(11\)}
    {\(12\)}
\end{problem}

\begin{answer}
\(11\)
\end{answer}

\begin{solution}
等比数列的首项为 \(a_{1}=1\)，所以 \(a_{n}=q^{n-1}\). 由题设，\(a_{1}a_{2}a_{3}a_{4}a_{5}=q^{0+1+2+3+4}=q^{10}\)，从而 \(q^{m-1}=q^{10}\). 等比数列的公比 \(q\ne0\)；若 \(m\ne11\)，则 \(q^{m-11}=1\)，这将推出 \(\lvert q\rvert=1\)，与题设矛盾. 所以 \(m=11\)，故选 C.
\end{solution}

\begin{problem}
一个长方体去掉一个小长方体，所得几何体的正（主）视图与侧（左）视图分别如图所示，则该几何体的俯视图为（\quad）

\texfigure[max width=0.72\linewidth]{img_repaint/2010/beijing_science/q3_condition_fixed.tex}

\choices
    {\choicetexfigure[width=0.15\paperwidth]{img_repaint/2010/beijing_science/q3_optA_fixed.tex}}
    {\choicetexfigure[width=0.15\paperwidth]{img_repaint/2010/beijing_science/q3_optB_fixed.tex}}
    {\choicetexfigure[width=0.15\paperwidth]{img_repaint/2010/beijing_science/q3_optC_fixed.tex}}
    {\choicetexfigure[width=0.15\paperwidth]{img_repaint/2010/beijing_science/q3_optD_fixed.tex}}
\end{problem}

\begin{answer}
C
\end{answer}

\begin{solution}
三视图满足“长对正、高平齐、宽相等”. 四个俯视图中的内折线分别表示去掉左上、右上、左下、右下的一个角. 主视图显示缺口位于左右方向的左侧；左视图显示缺口位于侧视图水平向的右侧. 按左视图与俯视图之间的深度方向反向对应关系，俯视图中的缺口应位于左下角. 因而 A 的左右位置不符，B 的左右位置和前后位置均不符，D 的左右位置不符，只有 C 同时还原题给的主视图与左视图，故选 C.
\end{solution}

\begin{problem}
8 名学生和 2 位老师站成一排合影，2 位老师不相邻的排法种数为（\quad）

\choices
    {\(\mathrm{A}_8^8\mathrm{A}_9^2\)}
    {\(\mathrm{A}_8^8\mathrm{C}_9^2\)}
    {\(\mathrm{A}_8^8\mathrm{A}_7^2\)}
    {\(\mathrm{A}_8^8\mathrm{C}_7^2\)}
\end{problem}

\begin{answer}
A
\end{answer}

\begin{solution}
先排 8 名学生，有 \(8!=\mathrm{A}_{8}^{8}\) 种排法. 学生排好后形成 9 个空位，包括两端的空位；两位老师不相邻，需从这 9 个空位中选出 2 个不同空位并安排两位不同老师，方法数为 \(\mathrm{A}_{9}^{2}\). 因此排法总数为 \(\mathrm{A}_{8}^{8}\mathrm{A}_{9}^{2}\)，故选 A.
\end{solution}

\begin{problem}
极坐标方程 \((\rho - 1)(\theta - \pi) = 0\) （\(\rho \geqslant 0\)） 表示的图形是（\quad）

\choices
    {两个圆}
    {两条直线}
    {一个圆和一条射线}
    {一条直线和一条射线}
\end{problem}

\begin{answer}
C
\end{answer}

\begin{solution}
由 \((\rho-1)(\theta-\pi)=0\)，得到两种情形：\(\rho=1\) 或 \(\theta=\pi\). 在极坐标中，\(\rho=1\) 表示以原点为圆心、半径为 1 的圆；在 \(\rho\geqslant0\) 的条件下，\(\theta=\pi\) 表示沿负 \(x\) 轴方向的一条射线. 因此图形是一个圆和一条射线，故选 C.
\end{solution}

\begin{problem}
\(\bs{a}\)，\(\bs{b}\) 为非零向量. “\(\bs{a} \perp \bs{b}\)”是“函数 \( f(x) = (x\bs{a} + \bs{b}) \cdot (x\bs{b} - \bs{a}) \) 为一次函数”的（\quad）

\choices
    {充分而不必要条件}
    {必要而不充分条件}
    {充分必要条件}
    {既不充分也不必要条件}
\end{problem}

\begin{answer}
B
\end{answer}

\begin{solution}
展开数量积得 \(f(x)=x^{2}\bs a\cdot\bs b+x(\bs b\cdot\bs b-\bs a\cdot\bs a)-\bs a\cdot\bs b\). 若 \(f(x)\) 是一次函数，二次项系数必须为零，所以 \(\bs a\cdot\bs b=0\)，即 \(\bs a\perp\bs b\)，因此前者是后者的必要条件. 但 \(\bs a\perp\bs b\) 时，若再有 \(\lvert\bs a\rvert=\lvert\bs b\rvert\)，则 \(f(x)\equiv0\)，不是一次函数，所以该条件不是充分条件. 故选 B.
\end{solution}

\begin{problem}
设不等式组 \( \left\{\begin{aligned}x+y-11&\geqslant0,\\ 3x-y+3&\geqslant0,\\ 5x-3y+9&\leqslant0\end{aligned}\right. \) 表示的平面区域为 \(D\). 若指数函数 \( y=a^{x} \) 的图象上存在区域 \(D\) 内的点，则 \(a\) 的取值范围是（\quad）

\choices
    {\((1,3]\)}
    {\([2,3]\)}
    {\((1,2]\)}
    {\([3, +\infty)\)}
\end{problem}

\begin{answer}
A
\end{answer}

\begin{solution}
由前两条边界 \(x+y=11\) 与 \(y=3x+3\) 相交于 \((2,9)\)，并由三条不等式可得区域 \(D\) 的横坐标满足 \(x\geqslant2\). 当 \(2\leqslant x\leqslant3\) 时，区域下界为 \(y=11-x\)；当 \(x\geqslant3\) 时，区域下界为 \(y=\frac{5}{3}x+3\)，上界始终为 \(y=3x+3\).

若 \(0<a\leqslant1\)，则 \(a^{x}\leqslant1\)，而 \(D\) 中的点满足 \(y\geqslant8\)，不可能相交. 若 \(a>3\)，对 \(x\geqslant2\) 有 \(a^{x}>3^{x}\geqslant3x+3\)，其中 \(3^{x}-3x-3\) 在 \(x=2\) 时为 0 且导数 \(3^{x}\ln3-3>0\)，故指数函数图象在区域上界之上，也不可能相交.

当 \(2<a\leqslant3\) 时，\(a^{2}\leqslant9\)，而 \(a^{3}>8\). 函数 \(a^{x}-(11-x)\) 在 \([2,3]\) 上连续，故在该区间内存在 \(x\) 使 \(a^{x}=11-x\)，对应点在 \(D\) 内. 当 \(1<a\leqslant2\) 时，\(a^{3}\leqslant8\)，而当 \(x\) 充分大时 \(a^{x}>\frac{5}{3}x+3\). 由连续性，在 \(x\geqslant3\) 内存在 \(x\) 使 \(a^{x}=\frac{5}{3}x+3\)，对应点也在 \(D\) 内. 因此 \(a\) 的取值范围为 \((1,3]\)，故选 A.
\end{solution}

\begin{problem}
正方体 \(ABCD - A_{1}B_{1}C_{1}D_{1}\) 的棱长为 2 ，动点 \(E\)，\(F\) 在棱 \(A_{1}B_{1}\) 上，动点 \(P\)，\(Q\) 分别在棱 \(AD\)，\(CD\) 上，若 \(EF = 1\)，\(A_{1}E = x\)，\(DQ = y\)，\(DP = z\)（\(x\)，\(y\)，\(z > 0\)），则四面体 \(PEFQ\) 的体积（\quad）

\choices
    {与 \(x\)，\(y\)，\(z\) 都有关}
    {与 \(x\) 有关，与 \(y\)，\(z\) 无关}
    {与 \(y\) 有关，与 \(x\)，\(z\) 无关}
    {与 \(z\) 有关，与 \(x\)，\(y\) 无关}
\end{problem}

\begin{answer}
D
\end{answer}

\begin{solution}
建立空间直角坐标系，令 \(A(0,0,0)\)，\(B(2,0,0)\)，\(C(2,2,0)\)，\(D(0,2,0)\)，\(A_{1}(0,0,2)\)，\(B_{1}(2,0,2)\). 则可设 \(E=(x,0,2)\)，且由于 \(EF=1\)，有 \(F=(x+\varepsilon,0,2)\)，其中 \(\varepsilon=\pm1\). 又 \(DP=z\)、\(DQ=y\)，故 \(P=(0,2-z,0)\)，\(Q=(y,2,0)\).

设四面体 \(PEFQ\) 的体积为 \(V\). 因为 \(\overrightarrow{PE}=(x,z-2,2)\)，\(\overrightarrow{PF}-\overrightarrow{PE}=\overrightarrow{EF}=(\varepsilon,0,0)\)，\(\overrightarrow{PQ}=(y,z,0)\)，所以 \(6V=\left\lvert\det(\overrightarrow{PE},\overrightarrow{EF},\overrightarrow{PQ})\right\rvert=2z\). 因此 \(V=\frac{z}{3}\)，只与 \(z\) 有关，与 \(x\)，\(y\) 无关，故选 D.
\end{solution}

\section{填空题}

\begin{problem}
在复平面内，复数 \( \frac{2\i}{1-\i} \) 对应的点的坐标为\fillinblank{}.
\end{problem}

\begin{answer}
\((-1,1)\)
\end{answer}

\begin{solution}
分子、分母同乘以分母的共轭复数，得 \(\frac{2\i}{1-\i}=\frac{2\i(1+\i)}{(1-\i)(1+\i)}=\frac{-2+2\i}{2}=-1+\i\). 因此该复数在复平面内对应的点为 \((-1,1)\).
\end{solution}

\begin{problem}
在 \(\triangle ABC\) 中，若 \(b = 1\)，\(c = \sqrt{3}\)，\(\angle C = \frac{2\pi}{3}\)，则 \(a = \)\fillinblank{}.
\end{problem}

\begin{answer}
\(1\)
\end{answer}

\begin{solution}
由余弦定理，\(c^{2}=a^{2}+b^{2}-2ab\cos C\). 代入 \(b=1\)、\(c=\sqrt{3}\)、\(C=\frac{2\pi}{3}\) 及 \(\cos\frac{2\pi}{3}=-\frac12\)，得 \(3=a^{2}+1+a\)，即 \(a^{2}+a-2=0\). 所以 \(a=1\) 或 \(a=-2\). 由于边长 \(a>0\)，故 \(a=1\).
\end{solution}

\begin{problem}
从某小学随机抽取 100 名同学，将他们的身高（单位：厘米）数据绘制成频率分布直方图（如图）. 由图中数据可知 \(a = \)\fillinblank{}. 若要从身高在 \([120, 130)\)，\([130, 140)\)，\([140, 150]\) 三组内的学生中，用分层抽样的方法选取 18 人参加一项活动，则从身高在 \([140, 150]\) 内的学生中选取的人数应为\fillinblank{}.

\texfigure[max width=0.72\linewidth]{img_repaint/2010/beijing_science/q11_hist.tex}
\end{problem}

\begin{answer}
\(0.03\)，\(3\)
\end{answer}

\begin{solution}
频率分布直方图中各小矩形的面积之和为 1，所以 \(10(0.005+0.035+a+0.020+0.010)=1\)，解得 \(a=0.030\). 三个指定组的总频率为 \(10(0.030+0.020+0.010)=0.6\)，其中身高在 \([140,150]\) 组的频率为 \(10\times0.010=0.1\). 分层抽样选取 18 人，故该组应选 \(18\times\frac{0.1}{0.6}=3\) 人.
\end{solution}

\begin{problem}
如图，\(\odot O\) 的弦 \(ED\)，\(CB\) 的延长线交于点 \(A\). 若 \(BD \perp AE\)，\(AB = 4\)，\(BC = 2\)，\(AD = 3\)，则 \(DE =\fillinblank{}\)； \(CE =\fillinblank{}\).

\texfigure[max width=0.52\linewidth]{img_repaint/2010/beijing_science/q12_circle.tex}
\end{problem}

\begin{answer}
\(5\)，\(2\sqrt{7}\)
\end{answer}

\begin{solution}
由图中 \(A\)，\(D\)，\(E\) 共线、\(A\)，\(B\)，\(C\) 共线，且 \(B\) 在 \(A\)，\(C\) 之间，得 \(AC=AB+BC=6\). 由相交弦的割线定理，\(AD\cdot AE=AB\cdot AC\)，所以 \(3AE=4\times6\)，得 \(AE=8\)，从而 \(DE=AE-AD=5\).

又因为 \(BD\perp AE\) 且 \(D\)，\(E\)，\(A\) 共线，所以 \(BD\perp DE\)，即 \(\angle BDE=90^\circ\). 因此 \(BE\) 是圆的直径，点 \(C\) 在圆上，故 \(\angle BCE=90^\circ\). 在直角三角形 \(ACE\) 中，\(CE^{2}=AE^{2}-AC^{2}=8^{2}-6^{2}=28\)，所以 \(CE=2\sqrt{7}\).
\end{solution}

\begin{problem}
已知双曲线 \(\frac{x^2}{a^2} - \frac{y^2}{b^2} = 1\) （\(a > 0\)，\(b > 0\)） 的离心率为 2，焦点与椭圆 \(\frac{x^2}{25} + \frac{y^2}{9} = 1\) 的焦点相同，那么双曲线的焦点坐标为\fillinblank{}；渐近线方程为\fillinblank{}.
\end{problem}

\begin{answer}
\((4,0)\)，\((-4,0)\)；\(y=\pm\sqrt{3}x\)
\end{answer}

\begin{solution}
椭圆 \(\frac{x^{2}}{25}+\frac{y^{2}}{9}=1\) 的焦半距为 \(c=\sqrt{25-9}=4\)，所以双曲线的焦点为 \((4,0)\)、\((-4,0)\). 双曲线离心率为 2，故 \(\frac{c}{a}=2\)，从而 \(a=2\). 再由 \(c^{2}=a^{2}+b^{2}\)，得 \(b^{2}=16-4=12\)，即 \(b=2\sqrt{3}\). 双曲线的渐近线为 \(y=\pm\frac{b}{a}x\)，所以方程为 \(y=\pm\sqrt{3}x\).
\end{solution}

\begin{problem}
如图放置的边长为 1 的正方形 \(PABC\) 沿 \(x\) 轴滚动.
设顶点 \(P(x, y)\) 的轨迹方程是 \(y = f(x)\)，则 \(f(x)\) 的最小正周期为\fillinblank{}；\(y = f(x)\) 在其两个相邻零点间的图象与 \(x\) 轴所围区域的面积为\fillinblank{}.

说明：“正方形 \(PABC\) 沿 \(x\) 轴滚动”包括沿 \(x\) 轴正方向和沿 \(x\) 轴负方向滚动. 沿 \(x\) 轴正方向滚动指的是先以顶点 \(A\) 为中心顺时针旋转，当顶点 \(B\) 落在 \(x\) 轴上时，再以顶点 \(B\) 为中心顺时针旋转，如此继续. 类似地，正方形 \(PABC\) 可以沿 \(x\) 轴负方向滚动.
\texfigure[width=0.24\linewidth]{img_repaint/2010/beijing_science/q14_fig1.tex}
\end{problem}

\begin{answer}
\(4\)，\(\pi+1\)
\end{answer}

\begin{solution}
正方形每经过一个顶点落到 \(x\) 轴上，支点向右移动一个边长 1；四个顶点依次落到 \(x\) 轴上后，正方形的姿态重复并向右平移 4，所以 \(f(x+4)=f(x)\). 在一个周期内，顶点 \(P\) 的相邻两次落地位置相距 4，且中点处不在 \(x\) 轴上，故最小正周期为 4.

从一次 \(P\) 落在 \(x\) 轴开始考察向右滚动，先以 \(P\) 为支点，直到 \(A\) 落地，此时 \(P\) 不移动；随后 \(P\) 绕 \(A\) 作半径为 1 的四分之一圆弧，绕 \(B\) 作半径为 \(\sqrt{2}\) 的圆心角为 \(\frac{\pi}{2}\) 的圆弧，再绕 \(C\) 作半径为 1 的四分之一圆弧，最后 \(P\) 再次落到 \(x\) 轴上. 两个半径为 1 的四分之一圆弧与 \(x\) 轴围成的面积各为 \(\frac{\pi}{4}\). 对中间圆弧，取其圆心 \(B\) 为原点，圆弧的方程为 \(y=\sqrt{2-u^{2}}\)，其中 \(-1\leqslant u\leqslant1\)，故其下方面积为 \(\displaystyle\int_{-1}^{1}\sqrt{2-u^{2}}\,\mathrm{d}u=1+\frac{\pi}{2}\). 因此相邻零点间的面积为 \(\frac{\pi}{4}+1+\frac{\pi}{2}+\frac{\pi}{4}=\pi+1\).
\end{solution}

\section{解答题}

\begin{problem}
已知函数 \( f(x) = 2\cos 2x + \sin^2 x - 4\cos x \).
\begin{enumerate}
\item 求 \( f\left(\frac{\pi}{3}\right) \) 的值；
\item 求 \( f(x) \) 的最大值和最小值.
\end{enumerate}
\end{problem}

\begin{solution}
\begin{enumerate}
\item
\[
f\left(\frac{\pi}{3}\right)=2\cos\frac{2\pi}{3}+\sin^{2}\frac{\pi}{3}-4\cos\frac{\pi}{3}=-1+\frac{3}{4}-2=-\frac{9}{4}
\]

\item 令 \(t=\cos x\)，则 \(-1\leqslant t\leqslant 1\)，并记 \(g(t)=f(x)\)，则
\[
g(t)=2\left(2t^{2}-1\right)+\left(1-t^{2}\right)-4t=3t^{2}-4t-1=3\left(t-\frac{2}{3}\right)^{2}-\frac{7}{3}
\]
因为 \(\frac{2}{3}\in[-1,1]\)，所以当 \(t=\frac{2}{3}\) 时取得最小值 \(-\frac{7}{3}\). 又因为 \(g(t)\) 是开口向上的二次函数，最大值在区间端点取得：\(g(-1)=6\)，\(g(1)=-2\). 因此，\(f(x)\) 的最大值为 \(6\)，最小值为 \(-\frac{7}{3}\).
\end{enumerate}
\end{solution}

\begin{problem}
如图，正方形 \(ABCD\) 和四边形 \(ACEF\) 所在的平面互相垂直，\(CE \perp AC\)，\(EF \parallel AC\)，\(AB = \sqrt{2}\)，\(CE = EF = 1\).

\begin{enumerate}
\item 求证：\(AF \parallel\) 平面 \(BDE\)；
\item 求证：\(CF \perp\) 平面 \(BDE\)；
\item 求二面角 \(A\)-\(BE\)-\(D\) 的大小.
\end{enumerate}

\texfigure[max width=0.56\linewidth]{img_repaint/2010/beijing_science/q16_solid.tex}
\end{problem}

\begin{solution}
设正方形所在平面为 \(z=0\)，建立空间直角坐标系，使 \(A(0,0,0)\)，\(B(\sqrt{2},0,0)\)，\(C(\sqrt{2},\sqrt{2},0)\)，\(D(0,\sqrt{2},0)\).
由题意，四边形 \(ACEF\) 所在平面与 \(z=0\) 垂直，且 \(CE\perp AC\)，\(CE=1\)，所以可取 \(E(\sqrt{2},\sqrt{2},1)\). 又因为 \(EF\parallel AC\)，\(EF=1\)，而 \(AC=2\)，故 \(F\left(\frac{1}{\sqrt{2}},\frac{1}{\sqrt{2}},1\right)\).

\begin{enumerate}
\item 取平面 \(BDE\) 的一个法向量
\[
\bs{n}=\overrightarrow{DB}\times\overrightarrow{DE}=(-\sqrt{2},-\sqrt{2},2)
\]
而 \(\overrightarrow{AF}=\left(\frac{1}{\sqrt{2}},\frac{1}{\sqrt{2}},1\right)\)，且 \(\overrightarrow{AF}\cdot\bs{n}=-1-1+2=0\).
又 \(\overrightarrow{BA}\cdot\bs{n}=2\neq0\)，故 \(A\) 不在平面 \(BDE\) 内. 所以 \(AF\) 的方向向量与平面 \(BDE\) 的法向量垂直，且 \(AF\) 不在平面 \(BDE\) 内，即 \(AF\parallel\) 平面 \(BDE\).

\item \(\overrightarrow{CF}=\left(-\frac{1}{\sqrt{2}},-\frac{1}{\sqrt{2}},1\right)\)，所以 \(2\overrightarrow{CF}=(-\sqrt{2},-\sqrt{2},2)=\bs{n}\).
因此 \(CF\) 的方向向量平行于平面 \(BDE\) 的法向量，从而 \(CF\perp\) 平面 \(BDE\).

\item 平面 \(ABE\) 的一个法向量和平面 \(DBE\) 的一个法向量分别可取 \(\bs{n}_{1}=\overrightarrow{AB}\times\overrightarrow{BE}=(0,-\sqrt{2},2)\)，\(\bs{n}_{2}=\bs{n}=(-\sqrt{2},-\sqrt{2},2)\).
于是 \(\bs{n}_{1}\cdot\bs{n}_{2}=6\)，\(|\bs{n}_{1}|=\sqrt{6}\)，\(|\bs{n}_{2}|=2\sqrt{2}\).
设二面角 \(A\)-\(BE\)-\(D\) 的平面角为 \(\theta\)，则取其锐角有
\[
\cos\theta=\frac{|\bs{n}_{1}\cdot\bs{n}_{2}|}{|\bs{n}_{1}||\bs{n}_{2}|}=\frac{6}{\sqrt{6}\cdot2\sqrt{2}}=\frac{\sqrt{3}}{2}
\]
所以 \(\theta=\frac{\pi}{6}\).
\end{enumerate}
\end{solution}

\begin{problem}
某同学参加 3 门课程的考试. 假设该同学第一门课程取得优秀成绩的概率为 \(\frac{4}{5}\)，第二、第三门课程取得优秀成绩的概率分别为 \(p\)，\(q\) （\(p > q\)）. 且不同课程是否取得优秀成绩相互独立. 记 \(\xi\) 为该生取得优秀成绩的课程数，其分布列为

\begin{center}
\begin{tabular}{c|cccc}
\(\xi\) & \(0\) & \(1\) & \(2\) & \(3\) \\
\hline
\(P\) & \(\frac{6}{125}\) & \(a\) & \(b\) & \(\frac{24}{125}\)
\end{tabular}
\end{center}

\begin{enumerate}
\item 求该生至少有 1 门课程取得优秀成绩的概率；
\item 求 \(p\)，\(q\) 的值；
\item 求数学期望 \(E(\xi)\).
\end{enumerate}
\end{problem}

\begin{solution}
\begin{enumerate}
\item 由分布列，\(P(\xi=0)=\frac{6}{125}\)，所以该生至少有 1 门课程取得优秀成绩的概率为
\[
P(\xi\geqslant1)=1-P(\xi=0)=1-\frac{6}{125}=\frac{119}{125}
\]

\item 由三门课程是否取得优秀成绩相互独立，有
\(P(\xi=0)=\frac{1}{5}(1-p)(1-q)=\frac{6}{125}\)，\(P(\xi=3)=\frac{4}{5}pq=\frac{24}{125}\).
两式分别化简得 \((1-p)(1-q)=\frac{6}{25}\) 和 \(pq=\frac{6}{25}\). 将前式展开并利用后式，得
\[
1-p-q+pq=pq\Longrightarrow p+q=1
\]
因此，\(p\)，\(q\) 是方程 \(t^{2}-(p+q)t+pq=0\)，即 \(t^{2}-t+\frac{6}{25}=0\) 的两个根. 解得
\[
t=\frac{1\pm\sqrt{1-\frac{24}{25}}}{2}=\frac{1\pm\frac{1}{5}}{2}
\]
所以 \(p=\frac{3}{5}\)，\(q=\frac{2}{5}\)（由 \(p>q\) 确定次序）.

\item
记第 \(i\) 门课程取得优秀成绩的指示变量为 \(I_i\)，则 \(\xi=I_1+I_2+I_3\)，且 \(E(I_1)=\frac{4}{5}\)，\(E(I_2)=p\)，\(E(I_3)=q\). 因此
\[
E(\xi)=\frac{4}{5}+p+q=\frac{4}{5}+1=\frac{9}{5}
\]
\end{enumerate}
\end{solution}

\begin{problem}
已知函数 \( f(x) = \ln\left(1 + x\right) - x + \frac{k}{2}x^2 \) （\(k \geqslant 0\)）.
\begin{enumerate}
\item 当 \(k=2\) 时，求曲线 \( y=f(x) \) 在点 \( (1,f(1)) \) 处的切线方程；
\item 求 \( f(x) \) 的单调区间.
\end{enumerate}
\end{problem}

\begin{solution}
\begin{enumerate}
\item 当 \(k=2\) 时，\(f(1)=\ln 2-1+1=\ln 2\)，且
\(f'(x)=\frac{1}{1+x}-1+kx\)，故 \(f'(1)=\frac{1}{2}-1+2=\frac{3}{2}\). 因此切线方程为 \(y-\ln 2=\frac{3}{2}(x-1)\)，即 \(3x-2y+2\ln 2-3=0\).

\item 函数的定义域为 \((-1,+\infty)\)，且
\(f'(x)=\frac{1}{1+x}-1+kx=\frac{x\left(kx+k-1\right)}{1+x}\).
当 \(k=0\) 时，\(f'(x)=-\frac{x}{1+x}\)，所以 \(f(x)\) 在 \((-1,0)\) 上单调递增，在 \((0,+\infty)\) 上单调递减.

当 \(0<k<1\) 时，令 \(r=\frac{1-k}{k}>0\). 由于定义域内 \(1+x>0\)，在 \((-1,0)\) 上有 \(x<0\) 且 \(kx+k-1<0\)，在 \((0,r)\) 上两因式异号，在 \((r,+\infty)\) 上两因式同为正. 因此，\(f(x)\) 在 \((-1,0)\) 和 \((r,+\infty)\) 上单调递增，在 \((0,r)\) 上单调递减.

当 \(k=1\) 时，\(f'(x)=\frac{x^{2}}{1+x}\geqslant0\)，且除 \(x=0\) 外为正，所以 \(f(x)\) 在 \((-1,+\infty)\) 上单调递增.

当 \(k>1\) 时，令 \(r=\frac{1-k}{k}\)，则 \(-1<r<0\). 在 \((-1,r)\) 上 \(x<0\) 且 \(kx+k-1<0\)，在 \((r,0)\) 上两因式异号，在 \((0,+\infty)\) 上两因式同为正. 因此，\(f(x)\) 在 \((-1,r)\) 和 \((0,+\infty)\) 上单调递增，在 \((r,0)\) 上单调递减. 综上，\(r=\frac{1-k}{k}\) 时得到题中各情形的单调区间.
\end{enumerate}
\end{solution}

\begin{problem}
在平面直角坐标系 \(xOy\) 中，点 \(B\) 与点 \(A(-1,1)\) 关于原点 \(O\) 对称，点 \(P\) 是动点，且直线 \(AP\) 与 \(BP\) 的斜率之积等于 \(-\frac{1}{3}\).
\begin{enumerate}
\item 求动点 \(P\) 的轨迹方程；
\item 设直线 \(AP\) 和 \(BP\) 分别与直线 \(x = 3\) 交于点 \(M\)，\(N\). 问：是否存在点 \(P\) 使得 \(\triangle PAB\) 与 \(\triangle PMN\) 的面积相等？若存在，求出点 \(P\) 的坐标；若不存在，说明理由.
\end{enumerate}
\end{problem}

\begin{solution}
\begin{enumerate}
\item 由中心对称性，\(B(1,-1)\). 设 \(P(x,y)\)，则直线 \(AP\) 和 \(BP\) 的斜率分别为 \(\frac{y-1}{x+1}\) 和 \(\frac{y+1}{x-1}\).
由于斜率存在，必须有 \(x\neq-1\) 且 \(x\neq1\). 由题意 \(\frac{y-1}{x+1}\cdot\frac{y+1}{x-1}=-\frac{1}{3}\)，化简得 \(3(y^{2}-1)=-(x^{2}-1)\).
所以动点 \(P\) 的轨迹方程为 \(x^{2}+3y^{2}=4\)，且 \(x\neq\pm1\).

\item 直线 \(AP\) 与 \(x=3\) 交于 \(M\)，直线 \(BP\) 与 \(x=3\) 交于 \(N\)，故
\(M\left(3,1+\frac{4(y-1)}{x+1}\right)=\left(3,\frac{x+4y-3}{x+1}\right)\)，\(N\left(3,-1+\frac{2(y+1)}{x-1}\right)=\left(3,\frac{2y-x+3}{x-1}\right)\).
于是
\[
y_N-y_M=\frac{2y-x+3}{x-1}-\frac{x+4y-3}{x+1}=-\frac{2(x-3)(x+y)}{x^{2}-1}
\]
直线 \(AB\) 的方程为 \(x+y=0\)，且轨迹与直线 \(AB\) 的交点只有 \((1,-1)\) 和 \((-1,1)\)，这两点因斜率不存在已被排除. 因此 \(x+y\neq0\).

三角形 \(PAB\) 的底边 \(AB\) 长为 \(2\sqrt{2}\)，点 \(P\) 到直线 \(AB\) 的距离为 \(\frac{|x+y|}{\sqrt{2}}\)，所以 \(S_{\triangle PAB}=|x+y|\). 又因为 \(MN\) 在直线 \(x=3\) 上，点 \(P\) 到该直线的距离为 \(|x-3|\)，故 \(S_{\triangle PMN}=\frac{1}{2}|x-3|\,|y_N-y_M|=\frac{(x-3)^{2}|x+y|}{|x^{2}-1|}\). 令两面积相等并约去正数 \(|x+y|\)，得到 \((x-3)^{2}=|x^{2}-1|\).
由轨迹方程知 \(|x|\leqslant2\). 当 \(|x|\geqslant1\) 时，
\[
(x-3)^{2}=x^{2}-1\Longrightarrow x^{2}-6x+9=x^{2}-1\Longrightarrow x=\frac{5}{3}
\]
当 \(|x|<1\) 时，\((x-3)^{2}=1-x^{2}\Longrightarrow 2x^{2}-6x+8=0\)，其判别式 \(36-64<0\)，无实根. 因此 \(x=\frac{5}{3}\). 代入轨迹方程，得 \(3y^{2}=4-\frac{25}{9}=\frac{11}{9}\Longrightarrow y=\pm\frac{\sqrt{33}}{9}\).
这两个点均满足 \(x\neq\pm1\)，所以存在所求点，坐标为 \(\left(\frac{5}{3},\frac{\sqrt{33}}{9}\right)\) 或 \(\left(\frac{5}{3},-\frac{\sqrt{33}}{9}\right)\).
\end{enumerate}
\end{solution}

\begin{problem}
已知集合 \( S_{n} = \{X \mid X = (x_{1}, x_{2}, \ldots, x_{n}), x_{i} \in \{0, 1\}, i = 1, 2, \cdots, n\} \) （\( n \geqslant 2 \)），对于 \( A = (a_{1}, a_{2}, \cdots, a_{n}) \)，\( B = (b_{1}, b_{2}, \cdots, b_{n}) \in S_{n} \)，定义 \(A\) 与 \(B\) 的差为 \( A - B = (|a_{1} - b_{1}|, |a_{2} - b_{2}|, \cdots, |a_{n} - b_{n}|) \)； \(A\) 与 \(B\) 之间的距离为 \( d(A, B) = \sum_{i=1}^{n} |a_{i} - b_{i}| \).
\begin{enumerate}
\item 证明： \(\forall A\)，\(B\)，\(C \in S_n\)，有 \(A - B \in S_n\)，且 \(d(A - C, B - C) = d(A, B)\)；
\item 证明： \(\forall A\)，\(B\)，\(C \in S_n\)，\(d(A, B)\)，\(d(A, C)\)，\(d(B, C)\) 三个数中至少有一个是偶数；
\item 设 \(P \subseteq S_n\)，\(P\) 中有 \(m (m \geqslant 2)\) 个元素，记 \(P\) 中所有两元素间距离的平均值为 \(\overline{d}(P)\). 证明： \(\overline{d}(P) \leqslant \frac{mn}{2(m - 1)}\).
\end{enumerate}
\end{problem}

\begin{solution}
\begin{enumerate}
\item 设 \(A=(a_{1},a_{2},\ldots,a_{n})\)，\(B=(b_{1},b_{2},\ldots,b_{n})\)，\(C=(c_{1},c_{2},\ldots,c_{n})\).
由于 \(a_i\)，\(b_i\in\{0,1\}\)，所以 \(|a_i-b_i|\) 只能取 \(0\) 或 \(1\). 因此
\[
A-B=(|a_{1}-b_{1}|,|a_{2}-b_{2}|,\ldots,|a_{n}-b_{n}|)\in S_n
\]
对每个 \(i\)，因为 \(a_i\)，\(b_i\)，\(c_i\in\{0,1\}\)，有
\[
\left||a_i-c_i|-|b_i-c_i|\right|=|a_i-b_i|
\]
事实上，当 \(c_i=0\) 时，左边为 \(|a_i-b_i|\)；当 \(c_i=1\) 时，左边为 \(|(1-a_i)-(1-b_i)|=|a_i-b_i|\). 于是
\[
d(A-C,B-C)=\sum_{i=1}^{n}\left||a_i-c_i|-|b_i-c_i|\right|=\sum_{i=1}^{n}|a_i-b_i|=d(A,B)
\]

\item 对固定的坐标 \(i\)，记 \(u_i=|a_i-b_i|\)，\(v_i=|a_i-c_i|\)，\(w_i=|b_i-c_i|\).
若 \(a_i=b_i=c_i\)，则 \((u_i,v_i,w_i)=(0,0,0)\)；否则，\(a_i\)，\(b_i\)，\(c_i\) 中有两个相等、一个不同，所以 \((u_i,v_i,w_i)\) 是 \((1,1,0)\) 的一个排列. 因此，无论哪种情况，\(u_i+v_i+w_i\) 都是偶数. 对 \(i=1\)，\(2\)，\(\ldots\)，\(n\) 求和，得到
\[
d(A,B)+d(A,C)+d(B,C)=\sum_{i=1}^{n}(u_i+v_i+w_i)
\]
为偶数. 若 \(d(A,B)\)，\(d(A,C)\)，\(d(B,C)\) 三个数全为奇数，则它们的和为奇数，矛盾. 因此三个数中至少有一个是偶数.

\item 设在集合 \(P\) 的 \(m\) 个元素中，第 \(i\) 个坐标等于 \(1\) 的元素有 \(k_i\) 个，则第 \(i\) 个坐标对所有无序元素对的距离贡献为
\[
k_i(m-k_i)
\]
因为只有一个元素该坐标为 \(1\)、另一个为 \(0\) 时，该坐标才贡献 \(1\). 又由
\[
k_i(m-k_i)=\frac{m^{2}}{4}-\left(k_i-\frac{m}{2}\right)^{2}\leqslant\frac{m^{2}}{4}
\]
可得所有无序元素对的距离总和满足
\[
\sum_{\{A,B\}\subseteq P}d(A,B)=\sum_{i=1}^{n}k_i(m-k_i)\leqslant\frac{nm^{2}}{4}
\]
而 \(P\) 中无序的两元素组共有 \(\mathrm C_m^2=\frac{m(m-1)}{2}\) 个，所以
\[
\overline{d}(P)=\frac{\displaystyle\sum_{\{A,B\}\subseteq P}d(A,B)}{\mathrm C_m^2}
\leqslant\frac{nm^{2}/4}{m(m-1)/2}=\frac{mn}{2(m-1)}
\]
命题得证.
\end{enumerate}
\end{solution}
